%% Loom fill-in notes — auto-generated from sources/pigeonhole.md
%% Spine: "Choose the boxes, then count."  (single-theorem; confidence 0.9)
%% Build: latexmk -xelatex main.tex   (or: bash compile.sh / ./compile.ps1)
\documentclass{loom}
\runningthread{pigeonhole — choose the boxes, then count}

% braced shorthands (unicode-math: subscripts/macros must be braced)
\newcommand{\ceil}[1]{\left\lceil #1 \right\rceil}
\newcommand{\floor}[1]{\left\lfloor #1 \right\rfloor}

\begin{document}

\loomcover{The Pigeonhole Principle}{Choose the boxes, then count.}{combinatorics, week 3}{2026}

% =====================================================================
\section*{The one idea, and how to read}
\warmth{0}\quad\whisper{One trivial counting fact; all the work is choosing what the boxes are.}

\begin{strand}
The principle itself is a one-line truth: crowd more objects than boxes, and some box is shared.
Every theorem and every puzzle below is \emph{the same fact} --- the only thing that changes,
and the only thing you actually have to invent, is \keyword{what counts as a box and what counts
as an object}. Read the statements; fill the keystone of each proof and the box-choice in each
application.
\end{strand}

\block{The master move (read it across)}
\begin{center}
\begin{tabularx}{\linewidth}{@{}>{\bfseries}l L@{}}
\toprule
objects ($n$) & the things you are forced to place \\
boxes ($k$)   & the labels you map them to --- \emph{choose these well} \\
the map       & a function objects $\to$ boxes; the cleverness lives here \\
the count     & if $n>k$, some box repeats; sharply, some box holds $\ge \fillin[1.2cm]$ \\
\bottomrule
\end{tabularx}
\end{center}
\trigger{you must guarantee a coincidence, but cannot say \emph{which} two things coincide.}
\recall{Fill the master count: how many in the fullest box, at least?}

\clearpage
\tableofcontents
\clearpage

% =====================================================================
\section{The principle}
\warmth{0}\quad\whisper{State the sharp version; the famous one-liner is just $n=k+1$.}

\begin{strand}
There is really one theorem here. The schoolbook version (``$n>k$ forces a shared box'') is the
special case $n=k+1$ of the sharp statement, so we lead with the sharp one and read the rest off it.
\end{strand}

\begin{theorem}[strong form; src \S3]
If $n$ objects are placed into $k$ boxes, some box contains at least
$\ceil{n/k}$ objects.
\end{theorem}

\begin{proof}
Suppose, for contradiction, that every box held at most $\ceil{n/k}-1$ objects.
\TODO{the keystone bound: add up over all $k$ boxes and reach $< n$}
That contradicts having placed all $n$ objects, so some box holds at least $\ceil{n/k}$.
\end{proof}

\recall{Why is the bound $\ceil{n/k}$ and not $\floor{n/k}$?}

\begin{corollary}[simple form; src \S2]
A function $f:A\to B$ between finite sets with $|A|>|B|$ cannot be injective.
\end{corollary}

\begin{definition}[injective, the clause that breaks; src \S2]
A map $f:A\to B$ is \keyword{injective} when it sends distinct inputs to distinct outputs, i.e.
$a\neq a' \implies \fillin[2.2cm]$. The proof's whole force is that injectivity would give
$|B|\ge|A|$.
\end{definition}

\begin{strand}
Take $n=k+1$ in the strong form: $\ceil{(k+1)/k}=2$, so some box holds two. That \emph{is} the
simple form. One theorem, read at two settings.
\end{strand}

% =====================================================================
\section{Choosing the boxes}
\warmth{0}\quad\whisper{Each puzzle is one sentence of setup: name the boxes, name the objects.}

\begin{strand}
Now the only skill. In each problem the answer is forced the instant you pick the right boxes;
the arithmetic is the master count from page one.
\end{strand}

\begin{example}[the matching pair; src \S4]
A drawer has black and brown socks; you draw in the dark. Boxes $=$ the two colours, objects $=$
socks drawn. With $3$ socks in $2$ colour-boxes, some colour has $\ceil{3/2}=2$: a pair. So you
need to draw \fillin[1cm].
\end{example}

\begin{example}[shared birth month; src \S4]
Among $13$ people, two share a birth month. Boxes $= 12$ months, objects $= 13$ people; since
$13>12$, two collide.
\end{example}

\begin{yourturn}
\textbf{Acquaintances ($R(3,3)=6$).} In any $6$ people, $3$ are mutual acquaintances or $3$ are
mutual strangers. Fix one person $P$. Of the other $5$, by pigeonhole at least $\ceil{5/2}=3$ fall
into one box --- ``acquaintances of $P$'' or ``strangers to $P$.'' Finish it: what are the boxes
here, and why do those $3$ settle the claim?
\workspace[3]
\end{yourturn}

% =====================================================================
\section{The non-obvious move}
\warmth{0}\quad\whisper{When boxes aren't handed to you, \emph{build a map} into a small label set.}

\begin{strand}
The deep applications hide a constructed function from the objects into a set provably smaller
than the domain. Define the labels well and the pigeonhole does the rest.
\end{strand}

\begin{theorem}[Erdős--Szekeres, special case; src \S5]
Any sequence of $n^2+1$ distinct reals contains a monotone subsequence of length $n+1$.
\end{theorem}

\begin{proof}
Label term $i$ with the pair $(u_i,d_i)$: the longest increasing run ending at $i$, and the
longest decreasing run ending at $i$. If no monotone run reaches $n+1$, every label lies in the
grid $\{1,\dots,n\}^2$, which has only $n^2$ cells.
\TODO{count: $n^2+1$ terms into $n^2$ label-cells, so two terms share a label}
Comparing two terms with the same label forces one of $u$ or $d$ to grow, a contradiction.
\end{proof}

\begin{yourturn}
The boxes here are the $n^2$ \emph{labels}, the objects are the $n^2+1$ \emph{terms}. Write the
one-line contradiction once two terms $i<j$ share a label $(u,d)$: what happens to $u_j$ or $d_j$?
\workspace[2]
\end{yourturn}

\loose{Generalise: $n$ objects in $k$ boxes give a box with $\ge\ceil{n/k}$. What is the dual
``some box has $\le\floor{n/k}$'' good for?}

\begin{strand}
\keyword{Takeaway.} The principle is a one-line counting fact; the craft is choosing the boxes ---
and, when they aren't given, \emph{building a map into a set too small to keep everyone apart}.
\end{strand}

\end{document}
